\section{Problem Setup}
\label{slab:sec:problem}

This section fixes the search space and the goal: the query language
\toolname operates over, the integrity constraints it can exploit, and
the problem it solves.

%\subsection{Query Space}
%\label{slab:sec:querylanguage}

\head{Query space}
The language, defined in Figure~\ref{slab:fig:querysyntax}, is an
expressive subset of SQL covering arbitrary nesting and join
structure.
%\slabcity currently supports a wide subset of SQL queries, including arbitrary nesting and join patterns. See Figure~\ref{slab:fig:querysyntax} for the formal language currently supported by \slabcity.



%\subsection{Integrity Constraints}
%\label{slab:sec:integrityconstraints}

\head{Integrity constraints}
Constraints both shrink the space of databases against which
equivalence must hold and, as Example~2 of
\cref{slab:sec:motivating-example} illustrated, create rewrite
opportunities. \slabcity supports the following kinds:
\begin{itemize}[leftmargin=*]
\item 
Primary and foreign keys. 
\item 
Comparisons within a row --- e.g.,  $\texttt{T.StartTime} < \texttt{T.EndTime}$. 
\item 
Implication constraints within a row. For example, $\texttt{T.EmailType} \neq \emph{spam} \rightarrow \texttt{T.action} = \texttt{NULL}$. 
%means if column \texttt{EmailType} is not \emph{spam} then column \texttt{action} is \texttt{NULL}. 
\item 
Whether or not a column can be \texttt{NULL}. 
\item 
Range constraints (for columns of integer or numeric types).
\item 
Enum types (e.g., column \texttt{Device} can draw from \{ \emph{S8}, \emph{iPhone} \}).
\end{itemize}

\begin{comment}
\begin{enumerate}
    \item primary keys
    \item foreign keys
\item comparative constraints within a row, e.g. \texttt{T.StartTime < T.EndTime} 
    \item implication constraints within the same row, e.g. \texttt{T.EmailType != "spam" -> T.action = NULL} means if the value of column \texttt{EmailType} is not \texttt{"spam"}, then the value of column \texttt{action} is \texttt{NULL}
        \item whether or not a column can be \texttt{NULL}
    \item range constraints (for int or numeric type columns)
    \item enum types, e.g. column \texttt{Device} can only take values from the set \texttt{\{"S8", "iPhone"\}}
    %\item \tofix{special constraints within a column, including 1) values of a column are incremental, 2) values of a column are consecutive.} \barzan{this one is a bit strange, what is the actual name of this kind of constraint in Oracle? auto increment is a feature not an IC afaik}     \xinyu{I think this one is very rare in the benchmarks. we do not have to include this one in the list.}
\end{enumerate}
\end{comment}


\head{Problem statement}
With the language and constraints in place, the problem is the
following.
%\subsection{Query Optimization Problem Statement}
%\label{slab:sec:problemstatement}

\begin{definition}
Given a query $\inputquery$ and an integrity constraint $\dataconstraint$, find a query $\query$ from our query language, such that the following two conditions hold: 
\begin{enumerate}[leftmargin=*]
\item 
$\query$ is semantically equivalent to $\inputquery$ with respect to $\dataconstraint$. 
That is, $\query$ and $\inputquery$ produce the same output on any database $\inputdb$ that meets the integrity constraint $\integrityconstraint$. 
\item 
$\query$ runs faster than $\inputquery$. Here, query performance is measured using \texttt{EXPLAIN} or the actual execution of the query.
\end{enumerate}
\label{slab:def:problem}
\end{definition}

Two aspects of the definition deserve a remark. Equivalence is
relative to the constraints, which is what licenses rewrites like
$\query_4$ in \cref{slab:sec:motivating-example}: they would be wrong
on arbitrary databases and are right on the databases that can
actually occur. And the definition asks for a faster query, not just
an equivalent one, which is why the system must surface multiple
equivalent candidates and rank them (\cref{slab:sec:perfranking}),
rather than stop at the first equivalence it can prove.



\begin{figure}[!t]
\footnotesize
\centering
\[
\arraycolsep=1.5pt\def\arraystretch{1.1}
\begin{array}{rrll}

\emph{Query} & 
\query & 
\grammareq & 
\emph{an input table} \\ 
& & \ \ \ | &  
\sqlselect \sqlspace \targetlist \sqlspace \sqlfrom \sqlspace \query \sqlspace \sqlwhere \sqlspace \sqlpredicate \\ 
& & \ \ \ | &  
\sqlselect \sqlspace \targetlist \sqlspace \sqlfrom \sqlspace \query \sqlspace \sqlgroupby \sqlspace \cols \sqlspace \sqlhaving \sqlspace \sqlpredicate \\ 
& & \ \ \ | &  
\sqlselect \sqlspace \targetlist \sqlspace \sqlfrom \sqlspace \query \sqlspace \sqlorderby \sqlspace \cols  \\ 
& & \ \ \ | &  
\query \sqlspace \sqlinnerjoin \sqlspace \query \sqlspace \sqlon \sqlspace \sqlpredicate \\ 
& & \ \ \ | &  
\query \sqlspace \sqlleftjoin \sqlspace \query \sqlspace \sqlon \sqlspace \sqlpredicate 
%& & \ \ \ | &  
%\query \sqlspace \sqlrightjoin \sqlspace \query \sqlspace \sqlon \sqlspace \sqlpredicate 
\\[2pt]

\emph{Target List} & 
\targetlist & 
\grammareq & 
[ \target \sqlspace \sqlas \sqlspace \colalias, \mydots, \target \sqlspace \sqlas \sqlspace \colalias ]
\\
& & \ \ \ | &  
\sqldistinct \sqlspace [ \target \sqlspace \sqlas \sqlspace \colalias, \mydots, \target \sqlspace \sqlas \sqlspace \colalias ] 
%\rui{add DISTINCT?}
\\[2pt]

\emph{Target} & 
\target & 
\grammareq & 
\col \ | \ \agg(\col) \ | \ \agg( \sqldistinct \sqlspace \col) 
%\rui{add DISTINCT?} 
\\
& & \ \ \ | &  
\agg(\col) \sqlspace \sqlover( \sqlpartitionby \sqlspace \cols \sqlspace \sqlorderby \sqlspace \cols )
\\ 
& & \ \ \ | &  
\windowfunc \sqlspace \sqlover( \sqlpartitionby \sqlspace \cols \sqlspace \sqlorderby \sqlspace \cols )
\\[2pt] 

\emph{Column List} & 
\cols & 
\grammareq & 
[\col, \mydots, \col]
\\[2pt]

\emph{Column} & 
\col & 
\grammareq & 
\colalias \ | \ \emph{a column from an input table} 
\\[2pt]

\emph{Condition} & 
\sqlpredicate & 
\grammareq & 
\sqlexpr \ \emph{op} \ \sqlexpr \ | \ \sqlpredicate \land \sqlpredicate \ | \ \sqlpredicate \lor \sqlpredicate 
\\[2pt] 

\emph{Expression} & 
\sqlexpr & 
\grammareq & 
\col \ | \ \agg(\col) \ | \ \agg(\sqldistinct \sqlspace \col) \ | \ \emph{const}
\\[2pt] 

\emph{Agg. Func.} & 
\agg & 
\grammareq & 
\sqlmax \ | \ \sqlmin \ | \ \sqlavg \ | \ \sqlsum \ | \ \sqlcount 
\\[2pt] 

\emph{Window Func.} & 
\windowfunc & 
\grammareq & 
\sqldenserank \ | \ \sqlrank 
\\[2pt] 

\end{array}
\]
\vspace{-15pt}
\caption{Syntax of our query language.}
\label{slab:fig:querysyntax}
\vspace{-20pt}
\end{figure}


